\subsection{Домашние задачи на вторую неделю}
\begin{problem}
    На окружности случайным образом выбираются две точки.
    Найдите среднее расстояние между ними.
    Приведите как минимум два варианта понимания фразы <<случайным образом>> в условии, приводящии к различным ответам.
\end{problem}
\begin{solution}
    \begin{itemize}
        \item
        \item
    \end{itemize}
\end{solution}
\begin{problem}
    Четыре человека $A, B, C, D$ становятся в очередь в случайном порядке.
    Найдите:
    \begin{enumerate}
        \item Условную вероятность того, что $A$~---~первый, если $B$~---~последний.
        \item Условную вероятность того, что $A$~---~первый, если $A$~---~не последний.
        \item Условную вероятность того, что $A$~---~первый, если $B$~---~не последний.
        \item Условную вероятность того, что $A$~---~первый, если $B$ стоит в очереди позже $A$.
        \item Условную вероятность того, что $A$ стоит в очереди раньше $B$, если известно, что $A$ раньше $C$.
    \end{enumerate}
\end{problem}
\begin{solution} \leavevmode
    \begin{enumerate}
        \item $\P(A = 1\mid B = 4) = \dfrac{2}{6} = \dfrac{1}{3}$.
        \item $\P(A = 1\mid A \neq 4) = \dfrac{1/3}{3/4} = \dfrac{4}{9}$.
        \item $\P(A = 1\mid B \neq 4) = \dfrac{\P(A = 1 \wedge B \neq 4)}{\P(B \neq 4)} = \dfrac{1/4 \cdot 4/6}{3/4} = \dfrac{2}{9}$.
        \item $\P(A = 1\mid A < B)\} = \dfrac{1/4}{1/2} = \dfrac{1}{2}$.
        \item $\P(A < B \mid A < C) = \dfrac{\P(A < B \wedge A < C)}{\P(A < C)} = \dfrac{2/6}{1/2} = 2/3$.
    \end{enumerate}
\end{solution}
\begin{problem}
    Пусть перед вами находится $N$ корзин с шарами, причём в $k$-ой лежит $k - 1$ красных и $N - k$ синих шаров.
    Вы выбираете равновероятно одну из этих урн и достаёте оттуда два шара подряд без возвращения.
    Пусть оказалось что первый шар красный.
    Найдите условную вероятность того, что второй шар -- синий.
\end{problem}
\begin{solution}
    В $k$-ой корзине лежит $N - 1$ шар из которых $k - 1$ красных и $N - k$ синих шаров.
    Тогда искомая вероятность равна
    \[
        \P(2 - blue \mid 1 - red) = \dfrac{\sum\limits_{k = 1}^N 1/N \cdot (k - 1)/(N - 1) \cdot (N - k)/(N - 2)}{\sum\limits_{k = 1}^N 1/N \cdot (k - 1)/(N - 1)} = \dfrac{1}{3}.
    \]
\end{solution}
\begin{problem}
    Рассмотрим абсолютно пьяного человека, выходящего из ресторана, расположенного на прямолинейной улице.
    Введём вдоль этой улицы систему координат так, чтобы ресторан был расположен в точке с нулевой координатой.
    Пусть за одну секунду человек перемещается с вероятностью $3/5$ в сторону увеличения координаты на $1$ метр, а с вероятностью $2/5$ в сторону уменьшения на то же расстояние.
    Предположим, с $-20$ метров (и далее в отрицательную сторону) начинается обрыв, а в точке с координатой $+20$ м расположен дом этого человека.
    Если он подойдет к обрыву, то непременно с него свалится, а если дойдёт до дома — уцелеет и больше никуда в этот вечер не пойдет.
    Найдите вероятность того, что он вернётся домой.
\end{problem}
\begin{solution}
    Пусть $p(x)$~---~вероятность достичь дома с отметки в $x$ метров.
    Из начальных условий $p(-20) = 0, p(20) = 1$.
    Очевидно, что \[p(x) = \dfrac{2}{5}p(x + 1) + \dfrac{3}{5}p(x - 1).\]
    Тогда, мы получаем что последовательность удовлетворяет характеристическому уравнению $2\lambda^2 - 5\lambda + 3 = 0$ и её корни~---~$\lambda = 1$ и $\lambda = 3/2$.
    Тогда общим решением является $p(x) = C_1 + C_2(3/2)^x$.
    Из начальных условий:
    \[
        \begin{cases}
            C_1 + C_2(3/2)^{-20} = 0; \\
            C_1 + C_2(3/2)^{20} = 1.
        \end{cases}
    \]
    То есть $C_2((3/2)^{20} - (3/2)^{-20}) = 1 \Ra C_2 = \dfrac{1}{3^{20}/2^{20} - 2^{20}/3^{20}}$,
    а $C_1 = -C_2(2/3)^{20}.$ Поскольку нас интересует $p(0)$, то ответом будет
    \[
        p(0) = C_1 + C_2 = \dfrac{3^{20}}{2^{20} + 3^{20}}.
    \]
\end{solution}
\begin{problem}
    Игральный кубик подбрасывают, пока не выпадет шесть очков.
    Найдите вероятность того, что это случилось на втором подбрасывании, если известно, что впервые шестёрка выпала после чётного числа бросаний.
\end{problem}
\begin{solution}
    \[
        \P(2 \mid "6") = \dfrac{\P("6" \mid 2)\P(2)}{\sum\limits_{k = 1}^{+\infty} \P("6"\mid 2k) \P(2k)} = \dfrac{1/6 \cdot 5/6}{\sum\limits_{k = 1}^{+\infty} 1/6 \cdot (5/6)^{2k - 1}} = \dfrac{11}{36}
    \]
\end{solution}
\begin{problem}
    Дано фиксированное число $n \geq 4$.
    Две точки выбираются наудачу из отрезка $[-n, n]$.
    Пусть $p$ и $q$~---~координаты этих точек.
    Найдите вероятность того, что квадратный трёхчлен $x^2 + px + q$ имеет вещественные корни.
\end{problem}
\begin{solution}
    У квадратного трёхчлена есть вещественные корни $\Longleftrightarrow p^2 - 4q \geq 0$.
    Тогда, из соображений геометрической вероятности, мы получаем следующее число:
    \[
        \P = 1 - \dfrac{2}{4n^2}\int_0^{2\sqrt{n}}n - \dfrac{x^2}{4}dx = 1 - \dfrac{2}{3\sqrt {n}}.
    \]
\end{solution}
\begin{problem}
    Три предприятия занимаются изготовлением одинаковых автомобилей, причём первое выпускает $25\%$ всей продукции, второе $35\%$, а третье -- $40\%$ соответственно.
    На первом предприятии $5\%$ брака, на втором $4\%$, а на третьем -- $2\%$.
    Приобретённый автомобиль оказался бракованным.
    Найдите вероятность того, что он был изготовлен на первом предприятии.
\end{problem}
\begin{solution}
    Пусть $b$~---~автомобиль бракованный, $g$~---~автомобиль хороший.
    Необходимо найти
    \[
        \P(1\mid b) = \dfrac{\P(b \mid 1)\P(1)}{\P(b \mid 1)\P(1) + \P(b \mid 2)\P(2) + \P(b \mid 3)\P(3)} = \dfrac{0.05 \cdot 0.25}{0.05\cdot 0.25 + 0.04 \cdot 0.35 + 0.02 \cdot 0.4} = \dfrac{25}{69}
    \]
\end{solution}
\begin{problem}
    Доказать, что $\P(A \cup B \cup C) = 1 - \P(\overline{A}|\overline{B}\cap \overline{C})\P(\overline{B}|\overline{C})\P(\overline{C})$.
\end{problem}
\begin{solution}
    Предположим, что мы имеем право брать условные вероятности по нужным нам событиям (чтобы не возникало сингулярностей).
    Тогда
    \[
        \P(\overline{A}|\overline{B} \cap \overline{C})\P(\overline{B}|\overline{C})\P(\overline{C}) = \P(\overline{A}\cap \overline{B}\cap \overline{C}) = 1 - \P(A \cup B \cup C).
    \]
\end{solution}

